\scanpage{39}\mspage{38}
% @continuation
also stable under passage to a subcountry. Finally, if $\unsure{C|\mathcal C}$ is finitistic, $\unsure{C'|C}$ too\unsure{?|:|.}
% @disj

\uuline{\deleted{Figures} disjoint} One says that two figures are disjoint if they are compatible, and if their intersection is empty (\unsure{no|only} common multistrata).
% @warn
% slanted
\margin{\uline{caution} For two figures (even elementary) to be disjoint, it \unsure{is} \uuline{necessary}, but \deleted{\unsure{not}} \inserted{necessarily} \unsure{sufficient|sufficient} that their shadows be disjoint (\unsure{that is} : \unsure{cp pairs|cp for} \unsure{simplify} for two loci !)}
% @warn2
\margin{\unsure{if|If} \deleted{\unsure{of}} they are compatible, \deleted{\unsure{and}} and \unsure{?}\par
\begin{tikzpicture}[x=1cm,y=1cm]\draw[->] (0,0) .. controls (.3,.35) and (.7,.25) .. (1,.05);\end{tikzpicture}}
% @objects
One says that \inserted{\uline{two multistrata}} $X,Y$ are \uline{disjoint} if $F_X,F_Y$ are, \inserted{i.e. if there is no \deleted{\unsure{$X$}\illegible{2}} $Z\in M$, with $Z\leq X,Y$}
% @equiv
% slanted
\margin{or what amounts to the same, if they are compatible and \deleted{the intersection} \deleted{of their} shadows are disjoint (\unsure{set-theoretically|set-theoretically})}
\margin{\begin{tikzpicture}[x=1cm,y=1cm]\draw (0,0) .. controls (1.4,.1) and (1.6,1.7) .. (1.3,2.4) .. controls (1.1,2.9) and (.3,3) .. (.05,2.75);\end{tikzpicture}}
% @all
% slanted
\margin{For $F,G$ to be disjoint, it is necessary and sufficient that $\forall X\in F$, $Y\in G$, $X$ and $Y$ be disjoint.}
\margin{\begin{tikzpicture}[x=1cm,y=1cm]\draw (1.3,3.2) .. controls (.7,3.15) and (.6,2.6) .. (.8,2.25) .. controls (1.05,1.7) and (1,1) .. (1,0);\end{tikzpicture}}
% @refine

\deleted{Let of} \deleted{two} figures $F,G$, one says that $F$ \uuline{refines} $G$ (and one writes $F\ll G$) if $\forall X\in F,\ \exists Y\in G$ such that $X\ll Y$. For two elementary figures $F_X,F_Y$ this therefore means \unsure{simply|simply} that $X\ll Y$.
% @prop

\deleted{If $F$ refines $G$} \uline{Proposition} For $F$ to be a subfigure of $G$, it is necessary and sufficient that $F$ refine $G$, and be compatible with $G$
\margin{\begin{tikzpicture}[x=1cm,y=1cm]\draw (0,0)--(0,1.15);\end{tikzpicture}}
% @proof

$\Longrightarrow$ trivial, to prove $\Longrightarrow$ one notes that \inserted{in the definition of refinement,} $X\ll Y$ \deleted{and} $X$ compatible with $Y$ \deleted{\unsure{and}} implies $X\leq Y$, hence $X\in G$ because of Fig 2.
